\section{Graph runtimes and opinionated harnesses}\label{sec:C22:graphs}
Graph runtimes make control flow an explicit application object. Nodes transform state, edges constrain possible continuations, and persistence attaches to defined execution boundaries. This can support inspection and fault recovery, but the graph is not automatically deterministic: nodes can contain model calls, time-dependent queries, or external effects. A graph's visual clarity therefore does not prove replay equivalence. The implementation must identify which values are recorded and which computations execute again.

LangGraph illustrates a runtime layer, while Deep Agents supplies a more opinionated harness over that ecosystem~\cite{executionproduct001,executionproduct006}. The latter's documented filesystem, subagent, and context-management interfaces are mechanisms to examine. Their existence does not establish that summaries preserve every constraint or that delegated work is cheaper. A useful experiment holds the underlying runtime fixed and separately varies offloading, compaction, and delegation. Otherwise a comparison between the two layers confounds control defaults with the runtime itself.

An opinionated harness reduces application effort by choosing defaults for prompts, tools, memory, and stopping. The benefit is faster construction of a coherent system. The scientific cost is that a high-level configuration can activate several mechanisms at once. Researchers should expose these defaults in their experimental configuration and inspect the code path used by the evaluated task. Disabling a component should not silently activate a fallback whose cost or behavior is materially different.

Event-driven workflow systems provide another control model. Events describe completed work or new requirements, and handlers may schedule further tasks. This is useful when external processes, human decisions, or delayed data determine progress. However, event duplication and ordering become part of agent semantics. A workflow that processes the same approval event twice can execute an action twice unless it deduplicates by stable event identity. These concerns persist even if every model call is perfectly accurate.

The discriminating experiment is not whether a sample application runs. It is whether equivalent state transitions produce equivalent committed effects under interruptions, retries, and schema evolution. Use a small deterministic tool fixture first, then introduce model nondeterminism as a separate factor. This ordering prevents a transport or persistence defect from being misdiagnosed as reasoning failure. It also establishes which guarantees survive when the framework delegates execution to a managed backend.

